\newcommand{\plotPGFfile}{true} % PGF olarak DATA dosyaları çizilsin mi çizilmesin mi? Çizilmez ise dosyadan sadece data okunan kısım atlanıyor, resmin çerçevesi oluşturuluyor...
\definecolor{dkgreen}{rgb}{0,0.6,0}
\definecolor{gray}{rgb}{0.5,0.5,0.5}
\definecolor{mauve}{rgb}{0.58,0,0.82}

\lstset{frame=tb,
	language=Java, %Uygulamanın yazıldığı programlama dili, (Örnek kullanım chapter4.tex'te verilmiştir.)
	aboveskip=3mm,
	belowskip=3mm,
	showstringspaces=false,
	columns=flexible,
	basicstyle={\small\ttfamily},
	numbers=none,
	numberstyle=\tiny\color{gray},
	keywordstyle=\color{blue},
	commentstyle=\color{dkgreen},
	stringstyle=\color{mauve},
	breaklines=true,
	breakatwhitespace=true,
	tabsize=3
}


\author{Halil}{Sarıkaya}{031021037}{031021037@std.izu.edu.tr}{Elektrik-Elektronik Mühendisliği}%1. Öğrenci ad soyad mail adresi bolum
%\authorIki{Ad}{Soyad}{ogrenciNo}{ad.soyad@std.edu.tr}{bolum}%2. Öğrenci ad soyad mail adresi bolum

%%%%
%\authorUc{Ad}{Soyad}{ogrenciNo}{ad.soyad@std.izu.edu.tr}{bolum}%3 Öğrenci ad soyad mail adresi bolum
%\authorDort{Ad}{Soyad}{ogrenciNo}{ad.soyad@std.izu.edu.tr}{bolum}%4. Öğrenci ad soyad mail adresi bolum
%\authorBes{Ad}{Soyad}{ogrenciNo}{ad.soyad@std.izu.edu.tr}{bolum}%4. Öğrenci ad soyad mail adresi bolum
\anabilimdali{Elektrik-Elektronik Mühendisliği}
\anabilimdaliUpper{ELEKTRİK-ELEKTRONİK MÜHENDİSLİĞİ}

\date{7 Haziran 2026}
\Yil{2025-2026}
\Donem{Güz}
\teziyoneten{Mohammed JOUDA, Dr. Öğr. Üyesi} % Danisman
\raporNo{Final}

\kisaltmalar{6DOF & 6 Serbestlik Derecesi \\
AHRS & Ataletsel Yönelim ve Referans Sistemi \\
CPU & Merkezi İşlem Birimi \\
DMA & Doğrudan Bellek Erişimi \\
DRDY & Veri Hazır Sinyali \\
EKF & Genişletilmiş Kalman Filtresi \\
EXTI & Dış Kesme Denetleyicisi \\
FPU & Kayan Nokta İşlem Birimi \\
FSM & Uçuş Durum Makinesi / Sonlu Durum Makinesi \\
GNSS & Küresel Konumlama Sistemi \\
GPIO & Genel Amaçlı Giriş/Çıkış \\
GPS & Küresel Konumlama Sistemi \\
GUI & Grafiksel Kullanıcı Arayüzü \\
HAL & Donanım Soyutlama Katmanı \\
I2C & Tümleşik Devreler Arası Haberleşme \\
IMU & Ataletsel Ölçüm Birimi \\
ISR & Kesme Servis Rutini \\
IWDG & Bağımsız Bekçi Zamanlayıcısı \\
LoRa & Uzun Menzilli Radyo Teknolojisi \\
LSI & Dahili Düşük Hızlı Saat Osilatörü \\
NVIC & İç İçe Kesme Denetleyicisi \\
PIL & Döngüde İşlemci Testi \\
RAM & Rastgele Erişimli Bellek \\
RTOS & Gerçek Zamanlı İşletim Sistemi \\
RX & Veri Alma / Alıcı Hattı \\
SIT & Sensör İzleme Testi \\
SPI & Seri Çevresel Arayüz \\
SRAM & Statik Rastgele Erişimli Bellek \\
SUT & Sentetik Uçuş Testi \\
TX & Veri Gönderme / Verici Hattı \\
UART & Evrensel Eşzamansız Alıcı-Verici \\
USART & Evrensel Eşzamanlı/Eşzamansız Alıcı-Verici}
\sembollistesi{$\theta$ & Yönelim / Eğim açısı \\
$\alpha$ & Tamamlayıcı filtre katsayısı \\
$\omega$ & Açısal hız \\
$\Delta t$ & Örnekleme periyodu / Zaman adımı \\
$\mathbf{e}$ & Hata vektörü \\
$\hat{\mathbf{g}}$ & Tahmini yerçekimi yönelim vektörü \\
$\mathbf{a}_{norm}$ & Normalize edilmiş ivme vektörü \\
$q$ & Yönelim kuaterniyonu / Kalman süreç gürültüsü katsayısı \\
$\dot{q}$ & Kuaterniyon türevi \\
$\otimes$ & Kuaterniyon çarpımı operatörü \\
$\omega_{gyro}$ & Jiroskop ölçüm vektörü \\
$\mathbf{b}$ & Jiroskop sapma (bias) vektörü \\
$K_p$ & Mahony filtresi oransal kazancı \\
$K_i$ & Mahony filtresi integral kazancı \\
$dt$ & Zaman adımı \\
$\hat{x}_k$ & Kalman filtresi durum vektörü \\
$F$ & Durum geçiş matrisi \\
$B$ & Giriş kontrol matrisi \\
$u_k$ & Kontrol girdisi vektörü \\
$P$ & Hata kovaryans matrisi \\
$Q$ & Süreç gürültüsü kovaryans matrisi \\
$R$ & Ölçüm gürültüsü kovaryans matrisi \\
$K_k$ & Kalman kazancı \\
$H$ & Gözlem matrisi \\
$z_k$ & Ölçüm vektörü \\
$r\_alt$ & Barometre ölçüm varyansı \\
$r\_acc$ & İvmeölçer ölçüm varyansı}




\title{MODEL ROKET UÇUŞ KONTROL BİLGİSAYARI} %TAMAMEN BÜYÜK HARF
\titleCase{Model Roket Uçuş Kontrol Bilgisayarı} %Sadece ilk Harfleri Büyük


